% Chapter 2

\chapter{LITERATURE SURVEY}\label{cha:lit}

Work on multi-objective UAV path planning comes from both graph search and evolutionary
optimisation~\cite{r11}. The main thing we asked was that once a planner holds a set of
good routes, how does it choose the one to fly, and does the payload get a say? Below we
show the papers which are nearest to that question and the few studies of what drone
transport does to medical payloads. The chapter ends with gaps we found and the
objectives we decided from them.

\section{PAPERS RELATED TO THE PROBLEM IDENTIFIED}\label{sec:lit-papers}

\subsection{Pareto-optimal particle swarm optimisation with entropy selection}

Chokkalingam and Ganesan~\cite{r1} propose a Pareto-optimal particle swarm optimiser
(POPSO) for UAV paths in changing surroundings, which has four parts. Grid-based and
random particles are mixed at the start which ensures that the swarm begins with varied
paths. Path length, energy, an inverse safety term, a penalty for moving obstacles and a
weather factor are the factors that determine the cost. Inertia weight and learning
coefficients shift during the run hence pushing exploration when progress stalls and
exploitation when it does not. Last, a hybrid Pareto-entropy step orders the archive by
Pareto rank and by an entropy value built from inter-solution distances, which keeps
the archive varied. On real terrain images with added obstacles it produced well-distributed fronts and
better path efficiency and safety than the methods it was compared with.

Two details matter for our work. The weather factor is
$W = \gamma\,(1 + v_\mathrm{wind}/v_\mathrm{max})$, which depends on wind speed only,
so a path flown into the wind and the same path flown with it are costed alike. And the
final route is not chosen by the entropy mechanism. For that the paper uses a separate
weighted sum of the normalised objectives, with weights ``typically determined
empirically based on mission priorities''. The same sum drives particle updates and
picks the one solution that gets deployed. So the last step is still
a hand-set weighted sum, and the payload does not enter it.

\subsection{Crowding-distance NSGA-II for urban UAV path planning}

Ren et al.~\cite{r3} plan UAV paths in an urban environment with two objectives,
distance and safety. They store the city as an octree, which splits 3-D space level by
level and uses little memory over large empty areas, and add a safety index map rating
each cell by its nearness to obstacles. A* paths seed the population and NSGA-II
evolves it. They show one case where the usual crowding distance favours an odd point,
and redefine crowding as the mean distance to the $k$ nearest solutions. The resulting
CDNSGA-II quickly finds well-spread, collision-free Pareto paths.

A weighted sum, they argue, needs its weights fixed beforehand, which is hard, so they
build the front first and choose from it later. The pick,
though, is made by
crowding distance, which measures how isolated a point is on the front and has no notion
of urgency, fragility or payload weight.

\subsection{Multi-strategy improved NSGA-II for 3-D trajectory planning}

Deng et al.~\cite{r11} go after two familiar faults of evolutionary planners in 3-D.
They get trapped in local optima, and fast convergence tends to cost population
diversity. Their MS-INSGA2 framework seeds the population with linearly
interpolated trajectories that pass a safety check, couples a dynamic guidance step
with crossover and mutation rates that adapt to the state of the population, and adds a
local search that refines the Pareto front and smooths the trajectories. Three
objectives are optimised: path length, collision risk from terrain and cylindrical
obstacles, and smoothness measured as the change in curvature between segments. On the
ZDT, DTLZ and WFG benchmark suites, with a Wilcoxon signed-rank test, the framework
outperformed its baselines, and on two 3-D planning tasks it beat eleven other
algorithms, lowering the average fitness by 27.87\% and 32.06\% and the best fitness by
6.91\% and 14.18\%.

The paper is about search quality and trajectory smoothness. Payload and delivery
context do not appear among the objectives, and wind is listed as future work rather
than modelled. As with the other evolutionary planners, the front is rebuilt for each
planning task.

\subsection{NSGA-II and MOPSO}

NSGA-II~\cite{r8} is still the reference evolutionary algorithm for several objectives.
It brought fast non-dominated sorting, elitism by merging parents and offspring, and the
crowding distance, which prefers solutions in sparse regions of the front. MOPSO~\cite{r6}
extended particle swarm optimisation with an external archive of non-dominated
solutions, leader selection on an adaptive grid and a mutation operator. Both are fast,
stochastic and widely used as baselines. Our MOPSO baseline adapts the second to graph
routes (Section~\ref{sec:m3}), and the crowding distance is one of the competing
selection rules that Tier~2 is compared against (Section~\ref{sec:m5}).

\subsection{Exact multi-objective A* search}

Mandow and P\'{e}rez de la Cruz~\cite{r17} analyse NAMOA*, a multi-objective extension
of A*. Instead of one cost per node, NAMOA* keeps a set of non-dominated cost vectors
(labels) per node and expands labels, not nodes. With an admissible heuristic it returns
the complete set of Pareto-optimal solution costs. The paper proves that, when the heuristic is
consistent, NAMOA* expands no more labels than any other admissible multi-objective
search algorithm given the same heuristic, and that a better-informed heuristic never
makes it expand more. NAMOA* is the starting point of our Tier~1 engine. Its
exactness is the reason we chose it over a population method: the front is the same on
every run, and interior points that a payload might need are never lost to chance.

\subsection{Payload-mass-aware trajectory planning}

Tsoutsouras et al.~\cite{r10} look at drones shared by several users, where packages of
different mass are dropped at different waypoints before their deadlines. They build a
power model in which propulsion power depends on both payload mass and flight speed,
and formulate trajectory planning as minimising deadline misses under energy,
waypoint and mass constraints, a problem close to the travelling salesman problem. A
greedy heuristic produces a first trajectory and an optimiser improves it, and the whole
computation runs on the drone, in under about 320\,ms on an ARM Cortex-A9.

This is the closest work to ours in that payload matters to the route. The way it
matters is different, though. Mass is a continuous physical parameter inside the
energy and speed model, and the planner returns one trajectory. There is no Pareto
front, and no discrete payload description such as fragility or urgency that chooses
between trade-offs.

\subsection{Quality indicators for Pareto fronts}

Zitzler and Thiele~\cite{r18} introduced the size of the dominated space, now called the
hypervolume (HV): the volume of objective space a front dominates up to a reference
point. It rewards both closeness and spread, and it is the main front-quality measure in
our benchmarks.

\subsection{Effect of drone transport on medical payloads}

Three studies measure what drone flight does to medical cargo. Oakey et al.~\cite{r13}
recorded vibration on fixed-wing and multi-copter drones and found levels up to three
times those of road transport, at higher frequencies, but insulin carried on those
flights showed no aggregation. Theobald et al.~\cite{r14} tried flying monoclonal
antibody infusions on sixty flights with vibration between 1.5 and 3\,g and found zero
change in particle size or aggregation. Stierlin et al.~\cite{r15} compared drone and car
transport of blood samples and found them comparable, with small differences in five
analytes.

So the payloads that were tested come unharmed. The studies
do not say what happens to a more fragile cargo, or whether a calmer corridor would
help, because none of them changed the route. We read this as a reason to treat
fragility as an operator's precaution and not as a proven damage model.

\subsection{Summary of the reviewed work}

The papers are shown in Table~\ref{tab:related-work}.

\begin{footnotesize}
\begin{longtable}{|P{0.5cm}|P{2.6cm}|P{4.1cm}|P{3.5cm}|P{3.6cm}|}
\caption{Summary of the reviewed work}\label{tab:related-work}\\
\hline
\textbf{No.} & \textbf{Author / Year} & \textbf{Method} & \textbf{Findings} & \textbf{Limitation for our problem}\\
\hline
\endfirsthead
\hline
\textbf{No.} & \textbf{Author / Year} & \textbf{Method} & \textbf{Findings} & \textbf{Limitation for our problem}\\
\hline
\endhead
1 & Chokkalingam and Ganesan (2026)~\cite{r1} & POPSO with hybrid initialisation, adaptive swarm parameters and Pareto-entropy archive & Well-distributed fronts; better path efficiency, safety and adaptability than compared methods & Final route chosen by a weighted sum with empirically set weights; wind enters as a speed-only factor; payload absent\\
\hline
2 & Ren et al. (2019)~\cite{r3} & CDNSGA-II over an octree city model with a safety index map; A* seeding & Spread-out, collision-free paths for distance and safety & Selection by crowding distance only; no mission context; no wind\\
\hline
3 & Deng et al. (2026)~\cite{r11} & MS-INSGA2 with safe interpolated seeding, adaptive operators and local search & Beats eleven algorithms on 3-D tasks; better benchmark indicators & No payload or delivery context; wind left to future work; front rebuilt per task\\
\hline
4 & Deb et al. (2002)~\cite{r8} & NSGA-II: fast non-dominated sorting, elitism, crowding distance & Better spread and convergence than PAES and SPEA on test problems & Crowding distance is purely geometric\\
\hline
5 & Coello Coello et al. (2004)~\cite{r6} & MOPSO with external repository, adaptive grid and mutation & Competitive fronts at low cost & Stochastic; may miss interior points of the front\\
\hline
6 & Mandow and P\'{e}rez de la Cruz (2010)~\cite{r17} & NAMOA*: label-based multi-objective A* & Exact front with admissible heuristics; efficiency results under consistent heuristics & Says nothing about which route to fly\\
\hline
7 & Tsoutsouras et al. (2020)~\cite{r10} & Mass- and speed-aware power model; heuristic plus optimiser for deadline-constrained multi-drop trajectories & Near-optimal trajectories on board in about 320\,ms & One trajectory, no front; mass is a physical parameter, not a selector\\
\hline
8 & Zitzler and Thiele (1999)~\cite{r18} & Comparative study of MOEAs; SPEA; dominated-space indicator & Hypervolume as a front-quality measure & Evaluation tool only\\
\hline
9 & Oakey et al. (2021); Theobald et al. (2023); Stierlin et al. (2024)~\cite{r13,r14,r15} & Vibration and sample-quality measurements on drone flights & Tested medicines and blood samples unaffected by ordinary flights & Route never varied; no evidence for or against route-based protection\\
\hline
\end{longtable}
\end{footnotesize}

\section{GAPS IDENTIFIED}\label{sec:gap}

Front diversity acts as the main criteria of design for two planners analyzed, but only
one of them uses it in choosing a route: Ren et al.~\cite{r3} choose a route using a
point selected directly from the front on the basis of crowding distance. Chokkalingam and
Ganesan~\cite{r1} use entropy only to keep the archive diverse during optimisation, and
fall back on a weighted sum, with weights set per mission, to make the final pick.
Deng et al.~\cite{r11} improve the front itself and leave the choice open.
Tsoutsouras et al.~\cite{r10} do not build a front at all: they put continuous payload
mass into the energy model and adapt speed as mass changes, but never choose between
competing trade-offs by payload.

So none of these works lets a payload property act as an input to a Pareto-point
selection step. This follows from how the mechanisms are defined: entropy and crowding
distance are functions of the front's geometry, and mass-aware adaptation is a function
of the physical model. Adding payload information would mean replacing the selection
rule outright, which is what our Tier~2 does.

We state the gaps as follows.

\begin{enumerate}
  \item[\textbf{Gap 1.}] No payload-aware route selection. Existing UAV
  multi-objective planners apply a fixed geometric criterion that ignores the payload,
  a hand-tuned weighted sum, or no selection mechanism at all.
  \item[\textbf{Gap 2.}] No reuse of a front between mission types. Every planner
  we read reruns the whole search when the mission changes. Nothing
  separates building the front, which depends on the region and the weather, from
  choosing a route, which depends on the payload.
  \item[\textbf{Gap 3.}] No algebraic bridge from payload to trade-off. We found
  no written, repeatable rule from a payload description to the trade-off it calls for
  (fragility to low risk, urgency to short distance, weight to low energy).
  \item[\textbf{Gap 4.}] Wind direction is not resolved per segment. When
  weather is modelled at all, it is one speed-only factor applied everywhere~\cite{r1}. A
  headwind and a tailwind on the same leg cost the same, and a stored front carries no
  record of the weather it was built for, so there is no rule for when it goes stale.
\end{enumerate}

\section{RESEARCH OBJECTIVES}\label{sec:objectives}

Each objective below answers one or more of these gaps.

\begin{enumerate}
  \item To cast UAV route planning as precomputing a Pareto front over a delivery
  airspace graph, and to justify exact search over stochastic methods.
  \item To build a 3-D airspace lattice from terrain, buildings, restricted zones, water
  and land use, with wind- and heading-dependent energy and risk on each edge, and a rule
  for when a stored front must be rebuilt.
  \item To define a payload metadata schema and a written formula turning mission
  context (fragility class, payload weight class, delivery urgency, priority) into
  weights for picking a Pareto point, with no expert tuning.
  \item To score the Tier~1 front against single-objective A* and MOPSO by
  hypervolume, inverted generational distance and spread.
  \item To gather evidence that Tier~2 picks routes that match the payload, including
  non-extreme routes for mixed payloads, and does so within one second.
  \item To export the selected route as an autopilot mission and verify it in a Gazebo
  and ArduPilot SITL simulation before hardware flights in Phase~II.
\end{enumerate}

Table~\ref{tab:gap-objective} links each gap to its objectives.

\begin{table}[!htbp]
\centering
\caption{Gaps identified and the research objectives that address them}\label{tab:gap-objective}
\small
\begin{tabular}{|P{3.6cm}|P{2.2cm}|P{8.4cm}|}
\hline
\textbf{Gap} & \textbf{Objectives} & \textbf{How the system addresses it}\\
\hline
Gap 1: no payload-aware selection & 3, 5 & Tier~2 derives weights from four payload tags and scores every route on the front; argmax gives the mission route\\
\hline
Gap 2: no front reuse & 1, 4 & Tier~1 builds an exact NAMOA* front once per region and weather window; a front store reuses it for every mission until it expires or the weather drifts\\
\hline
Gap 3: no algebraic bridge & 3, 5 & Weight formula $w_\mathrm{risk} \propto 2F$, $w_\mathrm{dist} \propto 3U$, $w_\mathrm{energy} \propto 2W$, with priority amplification; deterministic and checked for monotonicity over all 225 tag combinations\\
\hline
Gap 4: wind direction not resolved & 2, 6 & Headwind and crosswind computed per edge from its bearing; gust, turbulence and a go/no-go gate; front stored with its weather snapshot and refreshed on drift\\
\hline
\end{tabular}
\end{table}
